\section{One Crate per Board, and the Arrow Between Them}
\label{sec:bsp-one-crate}

The crate layout is not a matter of taste. It follows from one question: \emph{which of these things
can be owned separately?}

\subsection{The Two Schools}

\begin{table}[htbp]
    \centering
    \footnotesize
    \begin{tabular}{@{}p{0.24\textwidth}p{0.32\textwidth}p{0.36\textwidth}@{}}
        \toprule
        \textbf{Project} & \textbf{Layout} & \textbf{Where the wiring lives} \\
        \midrule
        \codeinline{stm32f4xx-hal} & one HAL per family, no BSP & in every example \\
        \codeinline{esp-hal} & one HAL, no BSP, chip by feature & in every example \\
        \codeinline{rp-hal} / \codeinline{rp-hal-boards} & HAL per chip + a BSP crate per board & in the BSP crate \\
        \bottomrule
    \end{tabular}
    \caption{Two schools. The difference is not sophistication: it is whether the board's wiring is a
    fact worth writing down once.}
    \label{tab:bsp-schools}
\end{table}

Both schools are legitimate. The first optimises for a HAL whose users all have a schematic in front of
them; the second optimises for a board whose wiring was expensive to establish. This project is firmly
in the second camp, for a reason that is in chapter~4's table of contradictions: the pin map of this
board was measured, and a measurement that is repeated in every example is a measurement that will be
repeated wrongly.

\subsection{The Four Crate Questions}

\begin{enumerate}
    \item \textbf{How many chips?} One HAL per chip or per family. This part gets one.
    \item \textbf{Do the boards stack?} If a carrier holds a module, one crate per piece, arrow carrier
          $\rightarrow$ module.
    \item \textbf{Are there variants?} A variant of the \emph{same} PCB is a feature. A different PCB is
          a crate.
    \item \textbf{Was the wiring obvious or was it conquered?} Obvious wiring can live in examples;
          conquered wiring belongs in a BSP with the measurement next to it.
\end{enumerate}

\subsection{What Ends Up in the Workspace}

\begin{minted}{text}
crates/
  efr32mg24-pac/            registers
  efr32mg24-hal/            drivers + pin and bus primitives
  xiao-mg24-bsp/            Module: led(), and with --features sense: imu(), microphone()
  xiao-expansion-base-bsp/  ExpansionBase: oled(), environment(), rtc(), buzzer(), button()
libraries/
  motion-detection/         domain policy: no hardware at all
\end{minted}

Two rules hold that shape together, and both are checks in this project's gate:

\begin{itemize}
    \item \textbf{Nothing in libraries/ may depend on the hardware stack.} A classifier
          that needs a HAL is a classifier that cannot be tested on a laptop, and the test suite is the
          only thing that makes it trustworthy.
    \item \textbf{Every crate that exposes hardware carries at least one example of its own level.} An
          API with no examples compiles, lints, and rots — this project had a whole module do exactly
          that for weeks.
\end{itemize}

\subsection{The Defect the Layout Prevented}

Before the split, the HAL itself contained \codeinline{board.rs} with the carrier's wiring — the OLED
bus, the button pin, the buzzer channel — because it was convenient. Two consequences followed
immediately:

\begin{enumerate}
    \item The HAL became unusable without the carrier's assumptions, so a user with a bare XIAO had to
          read carrier pins to blink an LED.
    \item Two BSPs existed at once: the HAL's board module and the new facade. Half the examples proved
          one, half the other, and the older one stopped compiling without anyone noticing.
\end{enumerate}

The fix was not better documentation. It was deleting the wiring from the HAL and making the
dependency direction a manifest fact.
